\section{Model: SVD Collaborative Filtering}
\label{sec:model}

\subsection{Singular Value Decomposition Foundation}

The core mathematical foundation of our model is the Singular Value Decomposition (SVD), a classical 
matrix factorization technique. For a matrix $R \in \mathbb{R}^{m \times n}$, the SVD decomposes it as:

\[
  R = U \Sigma V^T
\]

where:
\begin{itemize}
  \item $U \in \mathbb{R}^{m \times m}$ contains the left singular vectors (orthonormal basis of row space).
  \item $\Sigma \in \mathbb{R}^{m \times n}$ is a diagonal matrix of singular values, ordered in decreasing magnitude.
  \item $V \in \mathbb{R}^{n \times n}$ contains the right singular vectors (orthonormal basis of column space).
\end{itemize}

For recommendation, we use a truncated SVD with rank $k \ll \min(m, n)$:

\[
  R_k = U_k \Sigma_k V_k^T
\]

where $U_k$, $\Sigma_k$, $V_k$ use only the top $k$ singular values and corresponding vectors. This low-rank 
approximation captures the dominant structure while filtering noise and enabling generalization.

\subsection{Prediction Formula}

Given user $u$ and item $i$, the predicted rating is:

\[
  \hat{r}_{u,i} = U_{u,:} \Sigma_k (V_k^T)_{:,i} = \sum_{j=1}^{k} \sigma_j \cdot U_{u,j} \cdot V_{i,j}
\]

where:
\begin{itemize}
  \item $U_{u,:}$ is the $u$-th row of $U_k$ (user's latent factor representation).
  \item $(V_k^T)_{:,i}$ is the $i$-th column of $V_k^T$ (item's latent factor representation).
  \item $\sigma_j$ is the $j$-th singular value (importance of the $j$-th factor).
\end{itemize}

Intuitively, each latent factor captures a dimension of user preference and item characteristics. 
For example, one factor might represent ``action vs.~drama'', another ``popularity'', etc.~(though 
the factors are not labeled and are discovered automatically during decomposition).

\subsection{Surprise SVD Implementation}

The Surprise library implements SVD using stochastic gradient descent (SGD) rather than classical 
eigendecomposition. This is more practical for sparse, large-scale datasets and allows regularization.

\subsubsection{Training Algorithm}

The Surprise SVD minimizes:

\[
  L = \sum_{(u,i) \in \mathcal{R}} (r_{u,i} - \hat{r}_{u,i})^2 + \lambda \left(\|U\|_F^2 + \|V\|_F^2\right)
\]

using stochastic gradient descent. The updates for each observed rating $(u, i, r_{u,i})$ are:

\begin{align}
  U_{u.:} &\leftarrow U_{u,:} - \alpha \left(e_{u,i} \cdot V_{i,:} + \lambda \cdot U_{u,:}\right) \\
  V_{i,:} &\leftarrow V_{i,:} - \alpha \left(e_{u,i} \cdot U_{u,:} + \lambda \cdot V_{i,:}\right)
\end{align}

where:
\begin{itemize}
  \item $\alpha$ is the learning rate.
  \item $e_{u,i} = r_{u,i} - \hat{r}_{u,i}$ is the prediction error.
  \item $\lambda$ is the regularization parameter.
\end{itemize}

SGD has proven effective for collaborative filtering and converges reliably on the MovieLens dataset.

\subsubsection{Hyperparameters}

The Surprise SVD implementation is configured with the following hyperparameters:

\begin{table}[!h]
  \centering
  \small
  \begin{tabular}{llrp{7cm}}
    \toprule
    \textbf{Parameter} & \textbf{Symbol} & \textbf{Value} & \textbf{Rationale} \\
    \midrule
    Latent factors & $k$ & 100 & Balances expressiveness and overfitting. \\
    Learning rate & $\alpha$ & 0.005 & Provides stable convergence. \\
    Regularization & $\lambda$ & 0.02 & Prevents overfitting on sparse data. \\
    Epochs & $T$ & 20 & Sufficient for convergence on ML-100K. \\
    \bottomrule
  \end{tabular}
  \caption{SVD hyperparameter configuration.}
\end{table}

These values are chosen based on standard practice for the MovieLens 100K dataset and are held fixed 
for reproducibility in this laboratory setting.

\subsection{Handling Unknown Users and Items}

A critical design decision: how to handle predictions for users or items not in the training set?

The Surprise library provides a global mean rating $\mu$ as a fallback:

\[
  \hat{r}_{u,i} = \mu \quad \text{if } u \notin \mathcal{U} \text{ or } i \notin \mathcal{I}
\]

where $\mu = \frac{1}{|\mathcal{R}|} \sum_{(u,i) \in \mathcal{R}} r_{u,i}$.

This is a principled approach: for unknown entities, we predict the population average. In production, 
more sophisticated strategies (content-based features, user/item embeddings from alternative sources) 
could improve cold-start performance, but global mean provides a reasonable baseline.

\subsection{Rating Clamping}

Predictions from the model may occasionally fall outside the valid range $[1.0, 5.0]$ due to:
\begin{itemize}
  \item Extrapolation in factor space.
  \item Regularization effects pushing predictions away from extremes.
  \item Sparsity in user or item neighborhoods.
\end{itemize}

To ensure predictions always fall within the valid range, we apply post-processing:

\[
  \hat{r}_{u,i}^{\text{final}} = \max(1.0, \min(5.0, \hat{r}_{u,i}))
\]

This is simple and effective, though it discards probabilistic information. Production systems might 
instead return calibrated confidence intervals.

\subsection{Model Persistence and Serving}

The trained model is serialized using Python's \code{pickle} module and stored at \code{models/svd\_model.pkl}. 
Deserialization during API startup loads the model into memory, enabling $O(1)$ model lookup per prediction 
(after the initial factor computations).

Model versioning is handled through environment variables (\code{MODEL\_VERSION}), allowing A/B testing 
and rollback in production deployments.
